% !TeX root = closed-systems-from-comparison-completeness.tex
% ============================================================
\chapter*{Orientation to Part VI}
\phantomsection
\label{part:intro-electromagnetism}
% ============================================================
\begin{chapterorientation}
\orientationitem{Settled}{Part~V supplies connection-first geometry, the Hilbert-side scalar channel, and the stabilized quadratic carrier.}
\orientationitem{Task}{Part~VI derives phase curvature and then records the matter-sector scalar data that it supports.}
\orientationitem{Handoff}{The output is the charge lattice, the denominator--\(3\) composite-sector criterion, and the scalar/depth package used by numerical closure.}
\end{chapterorientation}

Part~VI treats electromagnetism as a phase-curvature consequence of the
closed transport stack.  The order is structural: phase reduction,
phase curvature, and charge quantization are developed from the
already-established carrier and connection data.

\Cref{ch:em} derives the \(U(1)\) phase sector, Maxwell equations in
the realized branch, and the charge-lattice consequences.  It also
records the composite-sector denominator result.  \Cref{ch:matter}
then identifies the scalar invariant and refinement-depth data assigned
to primitive matter sectors.  These data are the input to the numerical
closure part.
